\documentclass[10pt,a4paper]{amsart}
\usepackage{amsmath,amssymb,amsthm,graphicx}
% For commutative diagrams we use Xy-pic
\usepackage[all]{xy}

% We use 2cell for 2-commutative diagrams.
\xyoption{2cell}
% Theorem environments.
%
\theoremstyle{plain}
\newtheorem{theorem}[subsection]{Theorem}
\newtheorem{proposition}[subsection]{Proposition}
\newtheorem{lemma}[subsection]{Lemma}

\theoremstyle{definition}
\newtheorem{definition}[subsection]{Definition}
\newtheorem{example}[subsection]{Example}
\newtheorem{exercise}[subsection]{Exercise}
\newtheorem{situation}[subsection]{Situation}

\theoremstyle{remark}
\newtheorem{remark}[subsection]{Remark}
\newtheorem{remarks}[subsection]{Remarks}

\numberwithin{equation}{subsection}

% Macros
%
\def\lim{\mathop{\mathrm{lim}}\nolimits}
\def\colim{\mathop{\mathrm{colim}}\nolimits}
\def\Spec{\mathop{\mathrm{Spec}}}
\def\Hom{\mathop{\mathrm{Hom}}\nolimits}
\def\Ext{\mathop{\mathrm{Ext}}\nolimits}
\def\SheafHom{\mathop{\mathcal{H}\!\mathit{om}}\nolimits}
\def\SheafExt{\mathop{\mathcal{E}\!\mathit{xt}}\nolimits}
\def\Sch{\mathit{Sch}}
\def\Mor{\mathop{\mathrm{Mor}}\nolimits}
\def\Ob{\mathop{\mathrm{Ob}}\nolimits}
\def\Sh{\mathop{\mathit{Sh}}\nolimits}
\def\NL{\mathop{N\!L}\nolimits}
\def\CH{\mathop{\mathrm{CH}}\nolimits}
\def\proetale{{pro\text{-}\acute{e}tale}}
\def\etale{{\acute{e}tale}}
\def\QCoh{\mathit{QCoh}}
\def\Ker{\mathop{\mathrm{Ker}}}
\def\Im{\mathop{\mathrm{Im}}}
\def\Coker{\mathop{\mathrm{Coker}}}
\def\Coim{\mathop{\mathrm{Coim}}}

% Boxtimes
%
\DeclareMathSymbol{\boxtimes}{\mathbin}{AMSa}{"02}

%
% Macros for moduli stacks/spaces
%
\def\QCohstack{\mathcal{QC}\!\mathit{oh}}
\def\Cohstack{\mathcal{C}\!\mathit{oh}}
\def\Spacesstack{\mathcal{S}\!\mathit{paces}}
\def\Quotfunctor{\mathrm{Quot}}
\def\Hilbfunctor{\mathrm{Hilb}}
\def\Curvesstack{\mathcal{C}\!\mathit{urves}}
\def\Polarizedstack{\mathcal{P}\!\mathit{olarized}}
\def\Complexesstack{\mathcal{C}\!\mathit{omplexes}}
% \Pic is the operator that assigns to X its picard group, usage \Pic(X)
% \Picardstack_{X/B} denotes the Picard stack of X over B
% \Picardfunctor_{X/B} denotes the Picard functor of X over B
\def\Pic{\mathop{\mathrm{Pic}}\nolimits}
\def\Picardstack{\mathcal{P}\!\mathit{ic}}
\def\Picardfunctor{\mathrm{Pic}}
\def\Deformationcategory{\mathcal{D}\!\mathit{ef}}
\usepackage{cleveref}
\crefname{originalsection}{original section}{original sections}
\Crefname{originalsection}{Original section}{Original sections}
\makeatletter
\newlabel{section-gysin-higher-codimension}{{0FBI}{source}{}{sourcepointer-1}{}}
\newlabel{section-gysin-higher-codimension@cref}{{[originalsection][0][]0FBI}{[1][0][]source}{}{}{}}
\makeatother
\crefname{originallemma}{original lemma}{original lemmas}
\Crefname{originallemma}{Original lemma}{Original lemmas}
\makeatletter
\newlabel{lemma-gysin-easy}{{0FEA}{source}{}{sourcepointer-2}{}}
\newlabel{lemma-gysin-easy@cref}{{[originallemma][0][]0FEA}{[1][0][]source}{}{}{}}
\makeatother
\crefname{originallemma}{original lemma}{original lemmas}
\Crefname{originallemma}{Original lemma}{Original lemmas}
\makeatletter
\newlabel{algebra-lemma-CM-goes-up}{{045J}{source}{}{sourcepointer-3}{}}
\newlabel{algebra-lemma-CM-goes-up@cref}{{[originallemma][0][]045J}{[1][0][]source}{}{}{}}
\makeatother
\title{Intersection product of two integral subschemes of a smooth integral scheme equals the fundamental class of their intersection under a dimension and Cohen-Macaulay hypothesis}
\author{The Stacks Project Authors}
\date{stacks-project:chow.tex:local-snapshot-20261012; GNU FDL 1.2 or later}
\begin{document}
\maketitle
\noindent\textit{Editorial note.} All delivered mathematical text is the author's own, copied from the pinned Stacks Project chapter file chow.tex; no statement or proof step was rewritten, reordered or removed.\par\medskip
\noindent\textit{Editorial note.} Preamble provenance: the pinned chapter file opens with an input of the shared preamble and carries no preamble text of its own, so the project's single Stacks-wide preamble staged at sources/stacks-chow/preamble.tex is excerpted verbatim as two spans, lines 14-18 (Xy-pic setup) and lines 154-217 (the native theorem declarations and mathematical macros). Nothing was added, removed or redefined; the wrapper class and the amsmath/amssymb/amsthm/graphicx base packages are supplied by the assembler.\par\medskip
\noindent\textit{Editorial note.} The counted statement and proof contain no cross reference, no citation and no included file, so the delivered body is self contained.\par\medskip
\noindent\textit{Editorial note.} Canonical location of the counted proposition: Stacks tag 0FFD (sources/stacks-chow/pinned-original-source.tex line 14523).\par\medskip
\noindent\textit{Allowed standard prerequisites (editorial).}\begin{itemize}
\item The Gysin fundamental class identity for a virtual normal sheaf under dimension and Cohen-Macaulay hypotheses. Same chapter, Lemma labelled lemma-gysin-easy (source line 12128), invoked at proof lines 14547 and 14558-14559. Supplies the conclusion once the four hypotheses about the diagonal are verified; it is counted separately in this lane as S12-algebra-gysin-easy.
\item Higher codimension Gysin maps and the identification of the intersection product with a diagonal pullback. Same chapter, Section labelled section-gysin-higher-codimension (source line 11833), invoked at proof lines 14537-14540. Supplies that [Y] . [Z] = Delta\textasciicircum{}!([Y x Z]) with the Gysin map taken with respect to the locally free tangent sheaf.
\item Cohen-Macaulayness goes up along a flat local homomorphism with Cohen-Macaulay fibres. Algebra, Lemma labelled algebra-lemma-CM-goes-up (Stacks tag 045J), invoked at proof lines 14556-14557. Supplies that O\_\{Y x Z,xi\} is Cohen-Macaulay from the Cohen-Macaulayness of O\_\{Y,xi\} and the fibre O\_\{Z,xi\}.
\end{itemize}\medskip
\ifdefined\hypertarget\hypertarget{sourcepointer-1}{}\fi
\noindent\textit{Source reference.} Original section 0FBI: sources/stacks-chow/pinned-original-source.tex, source lines 11833--11833. Same-chapter pointer into the author environment lemma at source line 11833; invoked at proof line 14540 of the counted proof. It is retained only as a navigation pointer; every load-bearing step is inside the delivered span.\par\medskip
\ifdefined\hypertarget\hypertarget{sourcepointer-2}{}\fi
\noindent\textit{Source reference.} Original lemma 0FEA: sources/stacks-chow/pinned-original-source.tex, source lines 12128--12128. Same-chapter pointer into the author environment lemma at source line 12128; invoked at proof line 14547 of the counted proof. It is retained only as a navigation pointer; every load-bearing step is inside the delivered span.\par\medskip
\ifdefined\hypertarget\hypertarget{sourcepointer-3}{}\fi
\noindent\textit{Source reference.} Original lemma 045J: sources/stacks-chow/tags.txt, source lines 5438--5438. Cross-chapter standard prerequisite located in the project tag list staged in the local snapshot; invoked at proof line 14557 of the counted proof. It is retained only as a navigation pointer; every load-bearing step is inside the delivered span.\par\medskip
\begin{lemma}
\label{lemma-intersect-properly}
Let $k$ be a field. Let $X$ be an integral scheme smooth over $k$.
Let $Y, Z \subset X$ be integral closed subschemes. Set
$d = \dim(Y) + \dim(Z) - \dim(X)$. Assume
\begin{enumerate}
\item $\dim(Y \cap Z) \leq d$, and
\item $\mathcal{O}_{Y, \xi}$ and $\mathcal{O}_{Z, \xi}$
are Cohen-Macaulay for every $\xi \in Y \cap Z$ with
$\delta(\xi) = d$.
\end{enumerate}
Then $[Y] \cdot [Z] = [Y \cap Z]_d$ in $\CH_d(X)$.
\end{lemma}

\begin{proof}
Recall that $[Y] \cdot [Z] = \Delta^!([Y \times Z])$ where
$\Delta^! = c(\Delta : X \to X \times X, \mathcal{T}_{X/k})$
is a higher codimension Gysin map
(Section \ref{section-gysin-higher-codimension}) with
$\mathcal{T}_{X/k} = \SheafHom(\Omega_{X/k}, \mathcal{O}_X)$
locally free of rank $\dim(X)$. We have the equality of schemes
$$
Y \cap Z = X \times_{\Delta, (X \times X)} (Y \times Z)
$$
and $\dim(Y \times Z) = \dim(Y) + \dim(Z)$ and hence conditions 
(1), (2), and (3) of Lemma \ref{lemma-gysin-easy} hold.
Finally, if $\xi \in Y \cap Z$, then we have a flat local
homomorphism
$$
\mathcal{O}_{Y, \xi} \longrightarrow
\mathcal{O}_{Y \times Z, \xi}
$$
whose ``fibre'' is $\mathcal{O}_{Z, \xi}$. It follows that if both
$\mathcal{O}_{Y, \xi}$ and $\mathcal{O}_{Z, \xi}$
are Cohen-Macaulay, then so is $\mathcal{O}_{Y \times Z, \xi}$, see
Algebra, Lemma \ref{algebra-lemma-CM-goes-up}.
In this way we see that all the hypotheses of
Lemma \ref{lemma-gysin-easy} are satisfied and we conclude.
\end{proof}

\end{document}
